\section{Theoretical Analysis of Phase Transitions}
\label{sec:theory}

\subsection{Information-Theoretic Perspective}

The observed transitions at layers 3 and 8 can be understood through information-theoretic principles. Consider the mutual information $I(X; H^{(l)})$ between input $X$ and hidden representation $H^{(l)}$ at layer $l$. Our hypothesis is that phase transitions occur where the rate of information transformation changes:

\begin{equation}
\frac{d^2I(X; H^{(l)})}{dl^2} = 0 \text{ at } l \in \{3, 8\}
\end{equation}

This suggests layers 0-3 perform information compression (reducing redundancy), layers 3-8 perform structural transformation (syntactic parsing), and layers 8-12 perform semantic abstraction (concept formation).

\subsection{Linguistic Processing Hierarchy}

The three-phase structure aligns with established linguistic theory \cite{chomsky1957syntactic}:

\textbf{Phase 1 (Layers 0-3): Morphological Processing}
Early layers learn character and subword patterns, implementing robust tokenization. The 85\% recovery rate for character noise indicates redundant encoding that enables error correction through contextual patterns.

\textbf{Phase 2 (Layers 3-8): Syntactic Processing}
Middle layers construct hierarchical syntactic representations. The 78\% degradation under syntactic noise reflects the brittleness of tree-structured computations, where local errors propagate through dependency chains.

\textbf{Phase 3 (Layers 8-12): Semantic Processing}
Final layers build distributed semantic representations. The 67\% recovery rate suggests semantic compositionality provides alternative pathways for meaning reconstruction when syntax is corrupted.

\subsection{Computational Complexity Analysis}

We hypothesize that transition points minimize computational complexity while maximizing representational capacity. Let $C(l)$ denote computational cost and $R(l)$ denote representational power at layer $l$. The optimization problem:

\begin{equation}
\min_{\theta} \sum_{l=1}^{L} C(l) \quad \text{s.t.} \quad R(l) \geq R_{\text{min}}
\end{equation}

naturally produces phase boundaries where $\frac{dR(l)}{dC(l)}$ exhibits discontinuities, corresponding to our observed transitions.

\subsection{Gradient Flow Dynamics}

During training, gradient flow through transformer layers exhibits distinct patterns at phase boundaries. The gradient norm $\|\nabla_{\theta} \mathcal{L}\|$ shows peaks at layers 3 and 8, suggesting these layers learn qualitatively different transformations requiring larger parameter updates. This creates natural "checkpoints" in the computational graph where representations must be sufficiently stable to support downstream processing.